Using ideas from \Cref{sec:AppendixVL}, we prove a strong finite-page vanishing line on $\nu_{\BP}(\Ss^0)$. This line is not visible on the $\mathrm{E}_2$-page of the spectral sequence. The vanishing line will be used in \Cref{sec:Appendix} to provide the explicit numerical control over the function $\Gamma (k)$ required in \Cref{sec:Finale}.

\begin{cnv}
    In this section, we will fix a prime $p$ and implicitly $p$-localize all spectra. Furthermore, all synthetic spectra will be taken with respect to $\BP$.
\end{cnv}

\begin{thm}
  \label{thm:AN-vl}
  For $p\geq 3$, the $\BP$-synthetic sphere $\nu_{\BP} (\Ss^{0})$ has a strong finite-page vanishing line of slope $m$, intercept $c$ and torsion level $r$ where
  $$ m = \frac{1}{p^3 - p -1} \text{, }\ \ c = 2p^2 - 4p + 9 - \frac{2p^2+2p-10}{p^3 - p -1}\ \ \text{   and   }\ \  r = 2p^2-4p+2. $$
\end{thm}

\begin{rmk}
  The key content of \Cref{thm:AN-vl} is not the slope of the vanishing line,
  but rather the explicit values for the intercepts and torsion levels.\footnote{Indeed, an argument of Hopkins and Smith shows that the Devinatz--Hopkins--Smith nilpotence theorem is equivalent to the following statement:
  given any positive slope $\epsilon > 0$,
  the Adams-Novikov spectral sequence for $\Ss^0$ admits a vanishing line of slope $\epsilon$
  at some finite page (see \cite[Theorem 3.30]{Akhil} for a published account of this argument).
  By \Cref{prop:vanishing-line-Adams} this is equivalent to saying that $\nu_{\BP} \mathbb{S}^0$ has a finite page vanishing line of slope $\epsilon$ for every positive $\epsilon$.}
\end{rmk}

% The proof of \Cref{thm:AN-vl} consists of an analysis of the element $\beta_1$ in the homotopy groups of spheres. We start by choosing a synthetic lift of $\beta_1$ which will remain fixed for the remainder of the section. Choose a generator $\beta_1$ of $\pi_{qp-2}(\Ss^0)$ which is detected by $\beta_1$ on the 2-line of the Adams-Novikov spectral sequence $E_2$ term. 
% Then, by \Cref{thm:synthetic-Adams} we may choose an element $\wt{\beta}_1 \in \pi_{qp-2,qp}(\nu_{\BP}(\Ss^0))$ such that $\tau^{-1}\wt{\beta}_1 = \beta_1$, the image of $\wt{\beta}_1$ in $\pi_{qp-2,qp} (C\tau)$ is $\beta_1$ and $\wt{\beta}_1$ is not divisible by $\tau$.

\begin{ntn}
    We let $\wt{\beta}_1 \in \pi_{2p^2-2p-2,2p^2-2p} (\nu_{\BP} (\Ss^0))$ denote a synthetic lift of $\beta_1 \in \pi_{2p^2-2p-2} (\Ss^0)$ as in \Cref{thm:synthetic-Adams}(3).
\end{ntn}
The proof of \Cref{thm:AN-vl} consists of two main steps:
\begin{enumerate}
    \item We show that $C(\wt{\beta}_1)$ admits a strong vanishing line of slope $\frac{1}{p^3-p-1}$ and explicit intercept.
    \item Using the fact that $\beta_1$ is nilpotent topologically, we apply step (1) and the results of \Cref{sec:AppendixVL} to show that $\nu_{\BP} (\Ss^0)$ admits the desired strong finite-page vanishing line.
\end{enumerate}

%The proof of \Cref{thm:AN-vl} consists of an analysis of the element $\beta_1$ in the homotopy groups of spheres. We start by choosing a synthetic lift of $\beta_1$ which will remain fixed for the remainder of the section. By \Cref{thm:synthetic-Adams} we can choose an element $\wt{\beta_1} \in \pi_{qp-2,qp}(\nu_{\BP}(\Ss^0))$ such that $\tau^{-1}\wt{\beta_1} = \beta_1$ and the image of $\wt{\beta}_1$ in $\pi_{qp-2,qp} (C\tau)$ is the element $\beta_1$ (which detects $\beta_1$ in the Adams-Novikov spectral sequence).

%To prove \Cref{thm:AN-vl}, we will first prove that $C\wt{\beta}_1$ admits the desired vanishing line. Then we will use the nilpotence of $\beta_1$ to transfer this to a finite-page vanishing line in $\nu_{\BP} (\Ss^0)$. However, while $\beta_1$ is nilpotent as an element of $\pi_*(\Ss)$, it is not nilpotent on the $E_2$ term of the Adams-Novikov spectral sequence. As such $\wt{\beta}_1$ is not nilpotent as an element of the bigraded $\BP$-synthetic homotopy groups. Thus, $\Ss^{0,0}$ is not in the thick subcategory of $\BP$-synthetic spectra generated by $C\wt{\beta}_1$. However, from the fact that $\beta_1$ is nilpotent as an element of $\pi_*(\Ss)$ we can read off that $\wt{\beta}_1^N\tau^M = 0$ for $N,M \gg 0$. In terms of thick subcategories this implies that $\Ss^{0,0}$ is in the thick subcategory of $\BP$-synthetic spectra generated by $\{ C\wt{\beta}_1, C\tau \}$. Once we have shown that $C\wt{\beta}_1$ has the desired vanishing line we can use the techniques from \Cref{sec:AppendixVL} in order to finish the proof. We codify our argument in the following lemma.

%% We begin by providing a rough outline of the proof of Theorem \ref{thm:AN-vl}: 
%% \begin{enumerate}
%% \item Given an elemenent $b \in \pi_{*,*}(\Ss^{0,0})$ such that $b^N\tau^M = 0$ for $N,M \gg 0$ we prove a lemma which relates (strong) vanishing lines for $Cb$ and (strong) finite-page vanishing lines for $\Ss^{0,0}$.
%% \item Let $\wt{\beta}_1$ denote a generator of $\pi_{qp,qp-1}(\Ss^{0,0})$ in $\Syn_{\BP}$, then we show that $C\wt{\beta}_1$ has a strong vanishing line of the appropriate slope.
%% \item We apply the lemma described in step (1) to $b = \wt{\beta}_1$.
%% \end{enumerate}

%% ------------------- step 1 -------------------

Our proof of (1) will be based on the homological algebra of $P_*$-comodules, where $P_*$ is the polynomial part of the dual Steenrod algebra.

\begin{rec}
  Let $r$ denote the map of Hopf algebroids
  \[ (\BP_*, \BP_*\BP) \xrightarrow{r} (\F_p, P_*). \]
  Thinking in terms of the associated stacks we have a pullback/pushforward adjunction between the asscociated categories of sheaves
  \[ r^* : \mathrm{Stable}_{\BP_* \BP} \rightleftharpoons \mathrm{Stable}_{P_*} : r_*. \]
  Concretely, at the level of comodules $r^*$ sends a $\BP_*\BP$-comodule $M$ to the tensor product $\F_p \otimes_{\BP_*} M$ viewed as $P_*$-comodule.
\end{rec}

From \cite[Proposition 4.53]{Pstragowski} we know that the symmetric monoidal embedding \[\mathrm{Mod}_{C\tau} \to \mathrm{Stable}_{\BP_* \BP}\]
of \Cref{thm:tau-inv} is an equivalence.
As such, we will abuse notation by applying $r^*$ directly to $C\tau$-modules.

The reduction to $P_*$-comodules is carried out by the following lemma whose main input is an algebraic lemma of Krause \cite[Proposition 4.22]{AchimThesis}. For the statement of this lemma we use the notion of vanishing lines in Ext groups (which are a sort of bigraded homotopy groups) from Section 4 of Krause's work.

\begin{lem} \label{lemm:reduce-to-P}
  If $X$ is a compact and $\tau$-complete object of $\Syn_{\BP}$ such that
  $r^*(C\tau \otimes X)$ admits a vanishing line of slope $m$ and intercept $c$,
  then $X$ admits a strong vanishing line of the same slope and intercept.
\end{lem}

\begin{proof}
    By \Cref{lemm:vl-ctau} it suffices to show that $C\tau \otimes X \otimes \nu_{\BP} (A)$ admits a vanishing line of slope $m$ and intercept $c$ for all $A \in \Sp_{\geq 0}$. The vanishing statements we wish to prove are compatible with filtered colimits (as is $\nu_{\BP}$), therefore it suffices to restrict to the case where $A$ is finite.

  The symmetric monoidal equivalence of $\mathrm{Mod}_{C\tau}$ and $\mathrm{Stable}_{\BP_*\BP}$ provides us with a derived $\BP_*\BP$-comodule $\overline{X}$ associated to $C\tau \otimes X$ and equivalences 
    $$ \pi_{t-s,t}(C\tau \otimes X \otimes \nu_{\BP} (A)) \cong \Ext_{\BP_*\BP} ^{s,t} (\BP_*, \overline{X} \otimes_{\BP_*} \BP_*A). $$
    Let us now show that it suffices to address the case when $A = \Ss^0.$ The category of connective comodules over the Hopf algebroid $(\BP_*, \BP_*\BP)$ has enough projectives by the main result of \cite{SalchEnoughProj}, so we may fix a resolution $A_\bullet$ of $\BP_* A$ whose associated graded consists of positive shifts of $\BP_*$.
    We'll prove that the desired vanishing line already exists on the first page of the spectral sequence associated to the filtered object $\overline{X} \otimes_{\BP_*} A_\bullet$. By our choice of filtration this reduces to showing that $\Ext_{\BP_*\BP}^{s,t}(\BP_*, \overline{X})$ has the desired vanishing line, as we wanted.
  
    By our assumption that $X$ is compact, $\overline{X}$ is compact as an object of $\mathrm{Stable}_{\BP_* \BP}$. It therefore follows from the proof of \cite[Proposition 4.22]{AchimThesis} that $\overline{X}$ and $r^*\overline{X}$ admit the same vanishing lines, as desired.
\end{proof}

As a corollary to the above, we obtain the following well-known vanishing line. We will make use of it in our proof of \Cref{thm:AN-vl}.

\begin{prop}\label{prop:easy-ANSS-vl}
    Let $p$ be an odd prime. Then
  \label{prop:AN-vl}
  $\nu_{\BP}(\Ss^0)$ is $\tau$-complete and has a vanishing line of slope $m$ and intercept $c$ where
  $$ m = \frac{1}{p^2 - p -1}\ \ \text{   and   }\ \ c = 1 - \frac{2p-3}{p^2 - p -1}. $$
\end{prop}

\begin{proof}
    Since $\Ss^0$ is $\BP$-nilpotent complete, \Cref{lemm:Elocal-taucomplete} implies that $\nu_{\BP}(\Ss^0)$ is $\tau$-complete. Hence, by \Cref{lemm:reduce-to-P}, it suffices to show that the desired vanishing line is present in
  $$ \Ext_{P_*}^{s,t}(\F_p, \F_p). $$
    This vanishing line is already visible in the $\mathrm{E}_1$-page of the May spectral sequence for Ext over $P_*$.
\end{proof}

We now establish the desired vanishing line for $C(\wt{\beta}_1)$.

\begin{lem} \todo{It would probably be good from an expositional point of view to pull out the statement of \cite[Theorem 2.3.1]{PalmieriBook} from the proof of Lemma. I decided not to because of Palmieri's unfortunate notion of $CL$-spectrum. The other possibility would be trying to write our own theorem with explicit intercept in Achim's framework. I don't have time for this at the moment, though.}
  \label{lemm:cbeta-vl}
    The synthetic spectrum $C(\wt{\beta}_1)$ has a strong vanishing line of slope $m$ and intercept $c$, where
  $$ m = \frac{1}{p^3 - p - 1} \text{  and  } c = 8 - \frac{4p^2 - 11}{p^3 - p -1}. $$
\end{lem}

\begin{proof}
    We begin by noting that $C(\wt{\beta}_1)$ is $\tau$-complete by \Cref{prop:easy-ANSS-vl} and the fact that $\tau$-completeness is closed under finite colimits.
  Thus, by \Cref{lemm:reduce-to-P}, it suffices to show that 
    $ \Ext_{P_*}^{s,t}(\F_p, C(\beta_1)) $
    has the desired vanishing line, where $C(\beta_1)$ is the cofiber of the element $\beta_1 \in \Ext_{P_*} ^{2,2p^2-2p} (\F_p, \F_p)$ in $\mathrm{Stable}_{P_*}$.

    In \cite[Section 3]{Eva} Belmont shows that $C(\beta_1)$ satisfies the conditions of \cite[Theorem 2.3.1]{PalmieriBook} with Palmieri's paramter $d$ equal to $p^3 - p$. While this theorem is stated for $\A_*$ in \cite{PalmieriBook}, the proof carries over for $P_*$. Thus, we learn that
    $$ \Ext_{P_*}^{s,t}(\F_p, C(\beta_1)) = 0 \ \ \text{ for all } s > \frac{1}{d-1}(t-s + \alpha(d)) + 1, $$
  where $n$ is the minimal integer such that $2(p-1)p^n > d$ (in our case $n=2$) and
    $$ \alpha(d) := \left( \sum_{s+t \leq n \\ |\xi_t^{p^s}| \leq d} d + (p-1)|\xi_t^{p^s}| \right). $$
    The above calculation of the intercept is (the $P_*$ version of) \cite[Remark 2.3.4]{PalmieriBook}. We note that the $+1$ term at the end of the above inequality for $s$ comes from the $i_1$ term in \cite[Remark 2.3.4]{PalmieriBook}.
  We calculate that
  $\alpha(p^3 - p) = 7p^3 - 4p^2 - 7p + 4$ and thereby see that $C\wt{\beta}_1$ has the desired vanishing line.
\end{proof}

%% ------------------- step 2 -------------------

We now move on to step (2) of our proof of \Cref{thm:AN-vl}. Since $\beta_1$ is not nilpotent on the $\mathrm{E}_2$-page of the Adams-Novikov spectral sequence, the synthetic class $\wt{\beta}_1$ is not nilpotent. It follows that we cannot complete step (2) through a direct application of the genericity results of \Cref{sec:AppendixVL}. Instead, \Cref{thm:synthetic-Adams} will show that $\wt{\beta}_1 ^N \tau^M = 0$ for some large $N$ and $M$. In this situation, we have the following lemma:

\begin{lem}
  \label{lemm:syn-lines}
  Suppose that $X$ is a synthetic spectrum with a self map
  $b: \Sigma^{u,u+v}X \to X$ such that,
  \begin{itemize}
  \item $b^N\tau^M = 0$,
  \item $\Sigma^{-|b|} C(b)$ has a (strong) vanishing line of slope $m$ and intercept $c$, and
  \item $ \frac{v}{u} \geq m$.
  \end{itemize}
  Then $X$ has a (strong) finite-page vanishing line of slope $m$, intercept $c'$ and torsion level $M$, where
  $$ c' = c + \min(N(v-mu), M+m+1). $$
\end{lem}

\begin{proof}
%  We start with the following weaker argument:
%  it is clear that $X$ is in the thick subcategory generated by $\{C(b), C\tau\}$, thus by \Cref{thm:vl-are-generic} there exists a intercept $C$ and a torsion level $R$ such that $X$ has a (strong) vanishing line of slope $m$, intercept $C$ and torsion level $R$. The purpose of this lemma is to find explicit bounds on $C$ and $R$. 

  Consider the family of cofiber sequences
  $$ \Sigma^{-|b|} C(b) \to \Sigma^{-n|b|} C(b^n) \to \Sigma^{-n|b|} C(b^{n-1}) $$
  as $n$ varies. We will prove by induction that $\Sigma^{-n|b|} C(b^n)$ has a (strong) vanishing line of slope $m$ and intercept $c$. The base case is one of our hypotheses. In order to handle the induction step, we apply Lemma \ref{lemm:cof-line} to the cofiber sequence above. By assumption (and \Cref{lemm:suspension-line}), $\Sigma^{-n|b|} C(b^{n-1})$ has a (strong) vanishing line of slope $m$ and intercept $c + mu - v$. Thus, $\Sigma^{-n|b|} C(b^n)$ has a (strong) vanishing line of slope $m$ and intercept $\max(c, c + mu-v) = c$.
  
  Next, we apply Lemmas \ref{lemm:cof-line} and \ref{cor:ctau-vl} to the cofiber sequence 
  $$ C\tau^M \xrightarrow{f} \Sigma^{-N|b|} C(b^N\tau^M) \xrightarrow{g} \Sigma^{-N|b|} C(b^N) $$
  in order to conclude that $\Sigma^{-N|b|} C(b^N\tau^M)$ has a (strong) finite-page vanishing line of slope $m$, intercept $c$ and torsion level $M$. Finally, using the splitting
  $$ C(b^N\tau^M) \simeq X \oplus \Sigma^{(1,-M) + N|b|}X, $$
  we obtain the desired (strong) finite-page vanishing line. \qedhere
  
   %% left factor
   %% -Np, -N(p+q)
   %% m(x - Np) + c + Nq
   %% mx + c + N(q - mp)

   %% right factor
   %% 1, -M
   %% m(x+1) + c + M + 1
   %% mx + c + M + m + 1   
\end{proof}


%% \begin{proof}
%%   Our goal is to show that $C\wt{\beta}_1 \otimes \nu(A)$ has a vanishing line of slope $m$ and intercept $c$ for all $A \in \Sp_{\geq 0}$. 
  
%%   {\bf Step 1:}
%%   We the reduce to the statement that 
%%     $$ \Ext^{s,t}_{\BP_*\BP}(\BP_*, \BP_*/\beta_1) $$
%%   has the same vanishing line when plotted in the $(t-s,s)$-plane. Note here that $\BP_*/\beta_1$ refers to the cofiber of $\beta_1$ in the derived category of $\BP_* \BP$-comodules, and that the above Ext groups refer to its bigraded homotopy groups in this derived category.
  
%%     The vanishing statements we wish to prove are compatible with filtered colimits, (as is $\nu$) therefore it suffces to restrict to the case where $A$ is finite. By \Cref{lemm:vl-ctau} it suffices to show that $C\tau \otimes C\wt{\beta}_1 \otimes \nu_{\BP} (A)$ has the desired vanishing line. Here we use the fact hat $C\wt{\beta_1} \otimes \nu_{\BP} (A)$ is $\tau$-complete. Indeed, it follows from \Cref{rmk:E-complete} and \Cref{lemm:Elocal-taucomplete} that $\nu_{\BP} (A)$ is $\tau$-complete, so $C\wt{\beta}_1 \otimes \nu_{\BP} (A)$ is $\tau$-complete by the fact that taking cofibers preserves $\tau$-completeness. From \Cref{lemm:ctauE2} we find that
%%     $$ \pi_{t-s,t}(C\tau \otimes C\wt{\beta}_1 \otimes \nu_{\BP} (A)) \cong \Ext_{\BP_*\BP}(\BP_*, \BP_*/\beta_1 \otimes_{\BP_*} \BP_*A). $$
%%   In order to reduce to the case when $A = \Ss^0$ we will use an Atiyah-Hirzebruch spectral sequence argument. Fix a free resolution $X_\bullet$ of $\BP_*A$ as a $\BP_*\BP$-comodule such that the associated graded is just positive shifts of $\BP_*$. We'll prove that the desired vanishing line already exists on the first page of the spectral sequence associated to the filtered object $\BP_*/\beta_1 \otimes_{\BP_*} X_\bullet$. By our earlier identification of the associated graded this just mean showing that $\Ext_{\BP_*\BP}^{s,t}(\BP_*, \BP_*/\beta_1)$ has the desired vanishing line.

%%   {\bf Step 2:}
%%   The proof of Proposition 4.22 in \cite{AchimThesis} shows, in particular, that the vanishing line for the cofiber of $\beta_1$ as an endomorphism of the unit is the same over $BP_*BP$ and $P_*$, where $P_*$ refers to the polynomial part of the dual Steenrod algebra.
  
%%   {\bf Step 3:}
%%     In Section 3 of \cite{Eva} Belmont shows that the cofiber of $\beta_1$ (considerd as a derived $P_*$ comodule) satisfies the conditions of Theorem 2.3.1 from \cite{PalmieriBook} with $d = p^3 - p$. Note that this theorem is stated for $\A_*$ in \cite{PalmieriBook}, but the proof carries over easily for $P_*$. Thus, we learn that
%%   $$ \Ext_{P_*}^{s,t}(\F_p, \F_p/\beta_1) = 0 \ \ \text{ for all } s > \frac{1}{d-1}(t-s + \alpha(d)) + 1 $$
%%   where $n$ is the minimal integer such that $2(p-1)p^n > d$ (in our case $n=2$) and
%%     $$ \alpha(d) := \left( \sum_{s+t \leq n \\ |\xi_t^{p^s}| \leq d} d + (p-1)|\xi_t^{p^s}| \right). $$
%%     The above calculation of the intercept is (the $P_*$ version of) Remark 2.3.4 in \cite{PalmieriBook}. We note that the $+1$ term comes from the $i_1$ term in Remark 2.3.4
%%   A short calculation lets us conclude that
%%   $\alpha(p^3 - p) = 7p^3 - 4p^2 - 7p + 4$ and thereby that $C\wt{\beta}_1$ has the desired vanishing line.
%% \end{proof}

%% \begin{lem} \label{lemm:beta-1}
%%   $\pi_{qp-2,qp}(\Ss^{0,0}) \cong \F_p\{\tilde{\beta}_1\}$ where $\tilde{\beta}_1$ maps to $\beta_1 \in \pi_{qp-2}(\Ss^0)$ when we invert $\tau$ and $\beta_1 \in Ext_{\BP_*\BP}^{2,qp}(\BP_*,\BP_*)$ when we mod out by $\tau$.
%% \end{lem}

%% \begin{proof}
%%   This follows immediately from ???,
%%   the structure of the Adams-Novikov $E_2$ term through degree ??? as computed in CITE,
%%   and the fact that the first differential in the Adams-Novikov spectral sequence is in topological degree $2p^3 - 2p^2 - 3$.
%% \end{proof}

To apply this lemma to prove \Cref{thm:AN-vl}, we need to determine the constants that we have called $N$ and $M$ for $X = \nu_{\BP} (\Ss^0)$ and $b = \wt{\beta}_1$. By \Cref{thm:synthetic-Adams}, this comes down to the following lemma. 

\begin{lem}[Ravenel]\label{lemm:beta1-values}
    There are differentials
    \begin{align*}
        d_9 (\alpha_1 \beta_4) &= \beta_1 ^6 \text{ at } p=3 \text{ and} \\
        d_{33} (\gamma_3) &= \beta_1 ^{18} \text{ at } p=5
    \end{align*}
    in the Adams-Novikov spectral sequence. Moreover, we have 
    $\beta_1 ^{p^2-p+1} = 0$
    at any odd prime $p$. The Adams-Novikov differential with target $\beta_1 ^{p^2-p+1}$ has length at most
    $$2p^2 - 4p + 3. $$
\end{lem}

\begin{proof}
    The $3$-primary differential is part of \cite[Theorem 7.5.3]{GreenBook} and the $5$-primary differential is \cite[Theorem 7.6.1]{GreenBook}. 
    The general bound on the order of nilpotence of $\beta_1$ is proven shortly after the statement of Theorem 7.6.1 in \cite{GreenBook}, where Ravenel recounts a classical argument of Toda for this relation. Finally, the bound on the length of the differential follows from sparsity and the fact that there are no differentials off the $1$-line of the Adams-Novikov spectral sequence at odd primes.
\end{proof}

%% --------------- step 3 -----------------

\begin{proof}[Proof of Theorem \ref{thm:AN-vl}]
    In order to prove the theorem we apply \Cref{lemm:syn-lines} to $X = \nu_{\BP} (\Ss^{0})$ and $b = \wt{\beta}_1$. The remainder of the proof is just a matter of computing $m,c,u,v,N$ and $M$.

  The element $\wt{\beta}_1$ has bidegree $(2p^2-2p-2,2p^2-2p)$ so $u = 2p^2 -2p - 2$ and $v = 2$. By Lemmas \ref{lemm:cbeta-vl} and \ref{lemm:suspension-line} we know $\Sigma^{-|\wt{\beta}_1|}C\wt{\beta}_1$ has a strong vanishing line of slope $m = (p^3 - p - 1)^{-1}$ and intercept
  \begin{align*}
    c &= \left(8 - \frac{4p^2 - 11}{p^3 - p - 1} \right) - \left( 2 - \frac{2p^2 - 2p - 2}{p^3 - p - 1} \right) =  6 - \frac{2p^2+2p-9}{p^3 - p - 1}
  \end{align*}

    Suppose that there exists an $a$ in the $\mathrm{E}_{r+1}$-term of the Adams-Novikov spectral sequence such that $d_{r+1}(a) = \beta_1^N$. Then, by \Cref{thm:synthetic-Adams} there exists a $\wt{\beta_1^N}$ such that $\wt{\beta_1^N}\tau^{r} = 0$. A priori it may not be true that $\wt{\beta}_1^N = \wt{\beta_1^N}$, though we do know their difference maps to zero in $C\tau$ and is therefore divisible by $\tau$. In this case we can then use \Cref{prop:easy-ANSS-vl} to conclude that this ``difference divided by $\tau$'' is zero---seeing as it lives in a bigrading which is zero. To summarize, we learn that if $\beta_1^N$ is hit by a $d_{r+1}$-differential in the Adams-Novikov spectral sequence, then $\wt{\beta}_1^N \tau^r = 0$.
  
%  By \Cref{thm:synth-ASS}, it follows from the definition of $\wt{\beta}_1$ in \Cref{thm:synthetic-Adams} that $\wt{\beta}_1^N\tau^M$ must jump in $\nu \BP$-Adams filtration if there exists an $a$ such that $d_{M+1}(a) = \beta_1^N$. Moreover, it follows from \Cref{thm:synth-ASS} and \Cref{prop:easy-ANSS-vl}, interpreted as a statement about the $\mathrm{E}_2$-page of the Adams-Novikov spectral sequence for $\Ss^0$, that $\wt{\beta}_1 ^N \tau^M$ is already of minimal $\nu \BP$-Adams filtration in its bidegree. Therefore the existence of a differential $d_{M+1} (A) = \beta_1 ^N$ in fact implies that $\wt{\beta}_1^N \tau^M = 0$.

    We may therefore cite \Cref{lemm:beta1-values} to obtain the values of $N$ and $M$. We summarize the values we have computed in the following table:
  \renewcommand{\arraystretch}{1.6}
  \begin{center}
    \begin{tabular}{|c||c|c|c|c|c|}\hline
      prime  & $m$ & $u$ & $v$ & $N$ & $M$ \\\hline\hline
      $3$ & $\frac{1}{23}$ & $10$ & $2$ & $6$ & $8$ \\\hline
      $5$   & $\frac{1}{119}$ & $38$ & $2$ & $18$ & $32$ \\\hline
      $\geq 7$ & $\frac{1}{p^3 - p^2 - 1}$ & $2p^2-2p-2$ & $2$ & $p^2-p+1$ & $2p^2-4p+2$ \\\hline
    \end{tabular}
  \end{center}
  \renewcommand{\arraystretch}{1.0}
  At the prime $3$ the intercept is
  $$ 6 - \frac{15}{23} + \min \left(6\left( 2 - \frac{10}{23} \right), 9 + \frac{1}{23} \right) = 14 + \frac{9}{23}. $$
  At the prime $5$ the intercept is
  $$ 6 - \frac{51}{119} + \min \left(18 \left( 2 - \frac{38}{119} \right), 33 + \frac{1}{119} \right) < 38 + \frac{69}{119}. $$
  At primes $\geq 7$ the intercept is
  \begin{align*}
    6 - &\frac{2p^2+2p-9}{p^3 - p - 1} + \min \left((p^2 - p -1) \left( 2 - \frac{2p^2 - 2p - 2}{p^3 - p -1} \right), 2p^2 - 4p + 3 + \frac{1}{p^3 - p - 1} \right) \\
    &= 2p^2 - 4p + 9 - \frac{2p^2+2p-10}{p^3 - p^2 -1}.
  \end{align*}
  Note that the bound we write down for all primes is in fact equal to
  \[2p^2 - 4p + 9 + \frac{2p^2+2p-10}{p^3 - p^2 -1}. \qedhere\]
\end{proof}

%% \begin{rmk}\todo{AS:It's not super clear what phenomenon is being aimed at here.}
%%   Although the value of $N$ is smaller than $p^2 - p +1$ for small primes the value of $M$ is $2p^2 - 4p + 2$ for all primes. The authors wonder whether this a coincidence or part of a more general phenomenon.
%% \end{rmk}

%%% Local Variables:
%%% mode: latex
%%% TeX-master: "Main"
%%% eval: (visual-line-mode t)
%%% eval: (auto-fill-mode 0)
%%% End:
